\documentclass[11pt,a4paper]{amsart}

\usepackage[utf8]{inputenc}
\usepackage[T1]{fontenc}
\usepackage{amsmath,amssymb,amsthm}
\usepackage{mathtools}
\usepackage{hyperref}
\usepackage{microtype}
\usepackage{booktabs}
\usepackage{url}

\hypersetup{
  colorlinks=true,
  linkcolor=blue!60!black,
  citecolor=blue!60!black,
  urlcolor=blue!60!black,
}

% Theorem environments
\theoremstyle{plain}
\newtheorem{theorem}{Theorem}[section]
\newtheorem{proposition}[theorem]{Proposition}
\newtheorem{lemma}[theorem]{Lemma}
\newtheorem{corollary}[theorem]{Corollary}

\theoremstyle{definition}
\newtheorem{remark}[theorem]{Remark}

% Macros
\DeclareMathOperator{\ord}{ord}
\DeclareMathOperator{\Norm}{Norm}
\newcommand{\Z}{\mathbb{Z}}
\newcommand{\F}{\mathbb{F}}
\newcommand{\Q}{\mathbb{Q}}
\newcommand{\Phid}{\Phi_d(2)}
\newcommand{\wom}{\omega_3}
\newcommand{\Wp}{W_p}
\newcommand{\Wpm}{W_{p-2}}
\newcommand{\sqrttwo}{\sqrt{2}}

\title[Three BLS primality proofs for Wagstaff numbers]
      {Three Brillhart--Lehmer--Selfridge primality proofs \\ for Wagstaff numbers}

\author{Alexey Dolotov}
\address{}
\email{science@dolotov.com}

\date{\today}

\subjclass[2020]{Primary 11Y11; Secondary 11A51, 11B39}
\keywords{Wagstaff numbers, primality testing, Chebyshev polynomials,
  Pell sequences, Brillhart--Lehmer--Selfridge, APR-CL}

\begin{document}

\begin{abstract}
The Wagstaff numbers $W_p = (2^p + 1)/3$ for odd primes $p$ are the
natural $+1$ companions of the Mersenne numbers. Published primality
proofs for $W_p$ have relied on general-purpose elliptic-curve primality
proving (ECPP); Chebyshev/Lucas-type tests, while available as
compositeness criteria, have remained conjectural on the sufficiency
side. We give fully verified primality proofs of $W_{2617}$ (788
digits), $W_{10501}$ (3161 digits), and $W_{12391}$ (3730 digits),
independent of ECPP. The proofs apply the Brillhart--Lehmer--Selfridge
(BLS) $N-1$ framework~\cite{BLS75} to the cyclotomic decomposition
$2^{p-1} - 1 = \prod_{d \mid p-1} \Phi_d(2)$, with factorisations drawn
from the Cunningham project tables~\cite{Cunningham} and direct
ECM/$\rho/p{-}1$ computation. As an independent sanity check, the Chua
$N+1$ congruence $\omega_3^{(W_p + 1)/2} \equiv -1 \pmod{W_p}$ --- the
$a=3$ case of the Chua framework~\cite{Chua} with
$\omega_3 = 3 + 2\sqrt{2} \in \Z[\sqrt{2}]$, a necessary condition for
Wagstaff primality --- is verified at each $W_p$.

The method requires $p - 1$ to be sufficiently smooth that enough
cyclotomic factors $\Phi_d(2)$ have known factorisations; in the
current Cunningham tables~\cite{Cunningham} the exponents $p = 10501$
and $p = 12391$ are the only Wagstaff prime candidates above $2617$
meeting this ceiling. Every step is reproducible from the scripts
archived with this paper; all cofactor primalities are certified
unconditionally by APR-CL~\cite{APR,CohenLenstra}. The method does not
depend on the Chebyshev sufficiency conjecture; that conjecture is the
subject of a companion paper~\cite{PaperB}.
\end{abstract}

\maketitle

\section{Introduction}

\subsection{Wagstaff numbers and the ECPP era}

For an odd prime $p$, the \emph{Wagstaff number} is
\begin{equation*}
  W_p \;=\; \frac{2^p + 1}{3}.
\end{equation*}
The sequence $(W_p)$ is indexed in OEIS as A000979. All $W_p$ with
$p \leq 42737$ have been proved prime or composite; the exponents
$p \in \{5, 7, 11, 13, 17, 19, 23, 31, 43, 61, 79, 101, 127, 167,
191, 199, 313, 347, 701, 1709, 2617, 3539, 5807, 10501, 10691,
11279, 12391, 14479, 42737\}$ give the $29$ proved Wagstaff primes;
$7$ more probable primes are known up to the current computational
horizon. For exponents $p \geq 2617$, all published primality proofs
prior to this work used the elliptic-curve primality proving
algorithm of Atkin--Morain~\cite{AtkinMorain}.

The Mersenne counterparts $M_p = 2^p - 1$ admit the Lucas--Lehmer
test~\cite{Lehmer}, a deterministic standalone primality criterion. No
deterministic standalone test of comparable simplicity is known for
Wagstaff numbers. A natural Chebyshev/Lucas-type candidate --- the
single congruence
\begin{equation*}
  \wom^{(W_p + 1)/2} \;\equiv\; -1 \pmod{\Wp}, \qquad \wom = 3 + 2\sqrttwo,
\end{equation*}
in $\Z[\sqrttwo]$ (the $a = 3$ case of the Chua
framework~\cite{Chua}) --- is known to hold for every Wagstaff prime,
but its sufficiency remains conjectural; see the companion
paper~\cite{PaperB}.

In the absence of a deterministic standalone test, the
Brillhart--Lehmer--Selfridge (BLS) framework~\cite{BLS75} provides
unconditional primality proofs whenever $N \pm 1$ has a ``large
enough'' known factored portion. BLS is classical and venerable, but
the specific exponents $p = 2617, 10501, 12391$ have not previously
been put through BLS in the published literature; the ECPP proofs
already in the tables are of a completely different nature, and do not
depend on factorisations of $W_p \pm 1$.

\subsection{This paper}

We give three unconditional BLS primality proofs, one per exponent.
The proofs proceed by:
\begin{enumerate}
\item Establishing that Condition (II),
  $\wom^{(W_p+1)/2} \equiv -1 \pmod{W_p}$, holds for every Wagstaff
  prime $W_p$ (Proposition~\ref{prop:cond2}) --- this is the
  $a = 3$ case of the Chua framework~\cite{Chua} and follows from
  $W_p \equiv 3 \pmod{8}$.
\item Recording two derived criteria: a general $N+1$ criterion via
  $\wom$ (Theorem~\ref{thm:Nplus1}), useful when enough of $W_{p-2}$ is
  factored to give $F_+ > \sqrt{W_p}$; and its small-$p$ specialisation
  (Corollary~\ref{cor:small-p-specialisation}), which collapses to
  ``Condition (II) iff $W_p$ prime'' whenever $W_{p-2}$ itself is prime
  (Remark~\ref{rem:small-p}).
\item For the three target exponents $p \in \{2617, 10501, 12391\}$,
  applying the BLS $N - 1$ criterion (Theorem~\ref{thm:BLSN-1}) to the
  cyclotomic decomposition $2^{p-1} - 1 = \prod_{d \mid p - 1}
  \Phi_d(2)$, harvesting factored portions from the Cunningham project
  tables~\cite{Cunningham} plus direct ECM/$p{-}1$/Pollard-$\rho$
  computation. Every cofactor primality is certified by
  APR-CL~\cite{APR,CohenLenstra}.
\item Verifying Condition (II) independently at each of the three
  $W_p$ as a redundant primality witness --- consistent with
  Proposition~\ref{prop:cond2} and a useful debugging check on the
  modular-arithmetic implementation.
\end{enumerate}

The resulting proofs are short (Section~\ref{sec:proofs}), fully
unconditional, and cleanly separated from the Chebyshev sufficiency
conjecture (not used, not needed). All computations are reproducible
from the scripts described in Section~\ref{sec:comp}.

\subsection{Notation}

Throughout, $p \geq 5$ denotes an odd prime and $W_p = (2^p + 1)/3$.
For an integer or algebraic integer $x$ coprime to a modulus $m$,
$\ord_m(x)$ denotes the multiplicative order of $x$ modulo $m$;
$v_q(\cdot)$ denotes the $q$-adic valuation. We work in
$\Z[\sqrttwo]$, with fundamental unit $\alpha = 1 + \sqrttwo$ and
Chebyshev base $\wom = \alpha^2 = 3 + 2\sqrttwo$. The Pell sequences
$(U_n)$ and $(V_n)$ are defined by $\alpha^n = V_n/2 + U_n \sqrttwo$,
with the recurrence $X_{n+1} = 2 X_n + X_{n-1}$ and the identity
$V_n^2 - 8 U_n^2 = 4(-1)^n$. The Jacobi symbol is written $(a/n)$.

\section{Chebyshev bases and the Chua framework}

\subsection{The Chua identity}

For an integer $a \neq 0, \pm 1$, set $\omega_a = a + \sqrt{a^2 - 1}$,
so that $\omega_a \bar\omega_a = 1$ and $\omega_a^n + \omega_a^{-n} =
2 T_n(a)$ where $T_n$ is the $n$-th Chebyshev polynomial of the first
kind --- hence the name ``Chebyshev base''.

\begin{theorem}[Chua~\cite{Chua}; cf. Williams~{\cite[Ch.~4]{Williams}}]
\label{thm:chua}
Let $Q$ be an odd prime with $\gcd(a^2 - 1, Q) = 1$. Set
$D = a^2 - 1$, $\varepsilon = (D/Q)$, $\delta = (2(a+1)/Q)$. Then
\[
  \omega_a^{(Q - \varepsilon)/2} \;\equiv\; \delta \pmod{Q}
\]
in $\Z[\sqrt{D}]/(Q)$.
\end{theorem}

\begin{proof}
Define $\gamma = a + 1 + \sqrt{D}$. Then $\gamma^2 = 2(a+1) \cdot
\omega_a$. Compute $\gamma^Q$ in $\Z[\sqrt{D}]/(Q)$ by Frobenius
($\gamma^Q = a + 1 + \varepsilon\sqrt{D}$) and via the square
($\gamma^Q = \delta \cdot \omega_a^{(Q-1)/2} \cdot \gamma$). Equating
and dividing by $\gamma$ (a unit) yields $\omega_a^{(Q -
\varepsilon)/2} = \delta$.
\end{proof}

\subsection{Basic properties of Wagstaff numbers}

\begin{lemma}[Order of $2$]
\label{lem:order-2}
For every prime $p \geq 5$ and every prime $r \mid W_p$: $\ord_r(2) =
2p$ and $r \equiv 1 \pmod{2p}$.
\end{lemma}

\begin{proof}
Since $3 W_p = 2^p + 1$: $2^p \equiv -1 \pmod{r}$. Therefore
$\ord_r(2) \mid 2p$ but $\ord_r(2) \nmid p$. Since $p$ is prime,
$\ord_r(2) \in \{2, 2p\}$. If $\ord_r(2) = 2$: $r \mid 3$,
contradicting $\gcd(W_p, 3) = 1$ (since $p \geq 5$ is prime, $p \nmid
6 = \ord_9(2)$, so $v_3(2^p + 1) = 1$ and $3 \nmid W_p$). So
$\ord_r(2) = 2p$, and $2p \mid r - 1$ by Fermat.
\end{proof}

\begin{lemma}[mod-$8$ classification]
\label{lem:mod8}
For every odd prime $p \geq 3$:
\begin{enumerate}
\item[(i)]   $W_p \equiv 3 \pmod{8}$;
\item[(ii)]  every prime factor of $W_p$ satisfies $r \equiv 1$ or
             $3 \pmod{8}$;
\item[(iii)] the total multiplicity of prime factors $\equiv 3 \pmod 8$
             is odd.
\end{enumerate}
\end{lemma}

\begin{proof}
(i) Since $p$ is odd, $3 W_p = 2^p + 1 \equiv 1 \pmod 8$, so $W_p
\equiv 3 \pmod 8$. (ii) By Lemma~\ref{lem:order-2}, $v_2(\ord_r(2)) =
v_2(2p) = 1$. For $r \equiv 7 \pmod 8$: $(2/r) = 1$, so $\ord_r(2)
\mid (r-1)/2$ is odd, contradicting $v_2 = 1$. For $r \equiv 5 \pmod
8$: $v_2(\ord_r(2)) \leq v_2(r-1) = 2 \neq 1$. (iii) Write $W_p =
\prod r_i^{e_i}$. By (ii), each $r_i \equiv 1$ or $3 \pmod 8$. Modulo
$8$ the product is $\equiv 3^{\sum' e_i}$ where $\sum'$ runs over
$r_i \equiv 3$. From (i): $\sum' e_i$ is odd.
\end{proof}

\subsection{The base $a=3$}

\begin{proposition}[Base $a=3$ necessity]
\label{prop:cond2}
For every Wagstaff prime $W_p$ with $p \geq 5$,
\[
  \wom^{(W_p + 1)/2} \;\equiv\; -1 \pmod{W_p}, \qquad \wom = 3 + 2\sqrttwo.
\]
\end{proposition}

\begin{proof}
For $a = 3$: $D = a^2 - 1 = 8$ and $2(a + 1) = 8$, so $\varepsilon =
\delta = (8/W_p) = (2/W_p)^3 = (2/W_p)$. By
Lemma~\ref{lem:mod8}\,(i), $W_p \equiv 3 \pmod 8$, so the second
supplement to quadratic reciprocity gives $(2/W_p) = -1$; hence
$\varepsilon = \delta = -1$. Theorem~\ref{thm:chua} with $Q = W_p$
yields $\wom^{(W_p + 1)/2} \equiv -1 \pmod{W_p}$.
\end{proof}

\begin{remark}[Uniqueness]
\label{rem:uniqueness}
The base $a = 3$ is the unique integer $a \geq 2$ with $a^2 - 1$ a
power of $2$. Indeed, $(a-1)(a+1) = 2^k$ with $a - 1, a + 1$ two
consecutive even numbers forces $\gcd(a-1, a+1) = 2$, so both are
powers of $2$ differing by $2$ --- hence $\{a-1, a+1\} = \{2, 4\}$
and $a = 3$. Consequently, $a = 3$ gives the \emph{only} universal
$N + 1$ condition of the Chua form.
\end{remark}

The identity that powers our method:

\begin{lemma}[Fundamental identity]
\label{lem:fundamental}
For every $p \geq 5$:
\[
  W_p + 1 \;=\; 4\, W_{p-2}.
\]
\end{lemma}

\begin{proof}
$W_p + 1 = (2^p + 1)/3 + 1 = (2^p + 4)/3 = 4 \cdot (2^{p-2} + 1)/3 =
4 W_{p-2}$.
\end{proof}

\section{Primality criteria}

\subsection{Order lemma}

\begin{lemma}
\label{lem:order}
Let $N$ be odd with $N \equiv 3 \pmod 8$. Suppose
$\wom^{N+1} \equiv 1$ and $\wom^{(N+1)/2} \equiv -1 \pmod{N}$. Then
for every prime $r \mid N$: $v_2(\ord_r(\wom)) = v_2(N+1) = 2$.
\end{lemma}

\begin{proof}
The congruences descend from mod $N$ to mod $r$. Let $d = \ord_r(\wom)$.
From the first hypothesis: $d \mid N+1$. From the second: $d \nmid
(N+1)/2$. Hence $v_2(d) = v_2(N+1)$. Since $N \equiv 3 \pmod 8$: $N +
1 \equiv 4 \pmod 8$, so $v_2(N+1) = 2$.
\end{proof}

\subsection{The $N+1$ criterion}

\begin{theorem}[General $N+1$ criterion]
\label{thm:Nplus1}
Let $p \geq 5$ be prime, $N = W_p$. Suppose
Conditions (I) $\wom^{N+1} \equiv 1$ and (II) $\wom^{(N+1)/2} \equiv
-1 \pmod N$ hold. Let $S$ be a set of odd prime divisors of $W_{p-2}$
with $\wom^{(N+1)/q} - 1$ a unit modulo $N$ for each $q \in S$.
Define $F = 4 \prod_{q \in S} q^{v_q(W_{p-2})}$. If $F > \sqrt{N}$,
then $N$ is prime.
\end{theorem}

\begin{proof}
Suppose $N$ is composite with smallest prime factor $r \leq \sqrt N$.

\emph{Step 1 (Order constraints).}
By Lemma~\ref{lem:order}, $d := \ord_r(\wom)$ satisfies $d \mid N+1$
and $v_2(d) = 2$ (since $N + 1 = 4\,W_{p-2}$ with $W_{p-2}$ odd).

\emph{Step 2 (Full multiplicity: $F \mid d$).}
For each $q \in S$, the unit condition implies
$\wom^{(N+1)/q} \not\equiv 1 \pmod r$. Since $d \mid N+1$ but $d \nmid
(N+1)/q$: $v_q(d) = v_q(N+1) = v_q(W_{p-2})$. Together with
$v_2(d) = 2 = v_2(4)$: $F \mid d$.

\emph{Step 3 (Group-theoretic bound).}
Since $\Norm(\wom) = 9 - 8 = 1$ and $r$ is odd prime:
\begin{itemize}
\item If $(2/r) = -1$: $\Z[\sqrttwo]/(r) \cong \F_{r^2}$ and $\wom$
      lies in the norm-$1$ subgroup $\mu_{r+1}$, so $d \mid r + 1$.
\item If $(2/r) = +1$: $\Z[\sqrttwo]/(r) \cong \F_r \times \F_r$ and
      $d \mid r - 1$.
\end{itemize}

\emph{Case~$1$} ($d \mid r - 1$): $F \mid d \mid r-1$, so
$r \geq F + 1 > \sqrt N + 1$. Contradicts $r \leq \sqrt N$.

\emph{Case~$2$} ($d \mid r + 1$) --- boundary analysis.
From $F \mid d \mid r+1$: $r \geq F - 1$. Since $N \equiv 3 \pmod 8$,
$\sqrt N$ is irrational. With $F > \sqrt N$ integer: $F \geq \lfloor
\sqrt N \rfloor + 1$. Combined with $r \leq \lfloor \sqrt N \rfloor$:
$\lfloor \sqrt N \rfloor \geq r \geq F - 1 \geq \lfloor \sqrt N
\rfloor$, so $r = \lfloor \sqrt N \rfloor$ and $d = r + 1$.

Write $N = r \cdot s$ with $s \geq r$. From $(r+1) \mid (rs + 1)$ and
$r \equiv -1 \pmod{r+1}$: $s \equiv 1 \pmod{r+1}$. The smallest
candidate is $s = r + 2$; the next, $s = 2r + 3$, gives
$N = r(2r+3) = 2r^2 + 3r > r^2 + 2r + 1 = (r+1)^2$, contradicting
$r = \lfloor \sqrt N \rfloor$. So $N = r(r+2)$, $N + 1 = (r+1)^2$.
Since $N+1 = 4 W_{p-2}$: $W_{p-2} = ((r+1)/2)^2$. But $W_{p-2} \equiv
3 \pmod 8$ (Lemma~\ref{lem:mod8}), and a perfect square is $\equiv 0,
1, 4 \pmod 8$. Contradiction.
\end{proof}

\subsection{Coprimality and the $W_{p-2}$-prime case}

\begin{corollary}[Coprimality]
\label{cor:coprime}
Every prime divisor of $W_{p-2}$ is coprime to $W_p$.
\end{corollary}

\begin{proof}
$3 W_p - 12 W_{p-2} = (2^p + 1) - 4(2^{p-2} + 1) = -3$, so $W_p - 4
W_{p-2} = -1$.
\end{proof}

\begin{corollary}
\label{cor:small-p-specialisation}
If $W_{p-2}$ is prime, then $W_p$ is prime iff $\wom^{(W_p + 1)/2}
\equiv -1 \pmod{W_p}$.
\end{corollary}

\begin{proof}
\emph{Necessity}: Proposition~\ref{prop:cond2}. \emph{Sufficiency}: let
$r \mid N$. By Lemma~\ref{lem:order}, $d := \ord_r(\wom)$ has
$v_2(d) = 2$ and $d \mid N+1 = 4 W_{p-2}$. The divisors of $4 W_{p-2}$
with $v_2 = 2$ are $4$ and $4 W_{p-2}$. We exclude $d = 4$: $\wom^4
= 577 + 408\sqrttwo$, and $r \mid \gcd(576, 408) = 24$ contradicts
$\gcd(W_p, 24) = 1$. So $d = N + 1$. Since $d \mid r \pm 1$: $r \geq
N$, forcing $r = N$.
\end{proof}

\begin{remark}
\label{rem:small-p}
Corollary~\ref{cor:small-p-specialisation} applies directly to all
small-$p$ Wagstaff primes where $W_{p-2}$ is known prime: $p = 5$
(with $W_3 = 3$ prime), $p = 7$ ($W_5 = 11$), $p = 13$ ($W_{11} =
683$), $p = 19$ ($W_{17} = 43{,}691$), and ascends through the $29$
Wagstaff primes of §\ref{sec:appendix} as far as $p - 2$ also lies
in that list. The criterion in
Corollary~\ref{cor:small-p-specialisation} is a deterministic
primality test for $W_p$ in every such case; no factorisation beyond
confirming $W_{p-2}$ prime is required.
\end{remark}

\subsection{The BLS $N-1$ theorem}

For the three target exponents, $W_{p-2}$ is not known to have enough
factored content to satisfy $F_+ > \sqrt{W_p}$ in
Theorem~\ref{thm:Nplus1}; the $N + 1$ factored portion is $F_+ = 4$ at
best, because full factorisations of $W_{p-2}$ are unavailable at
these sizes. The proofs in Section~\ref{sec:proofs} therefore work on
the $N - 1$ side using the classical BLS theorem, which we recall:

\begin{theorem}[BLS $N-1$, {\cite[Theorem~5; strengthened form, Theorem~11]{BLS75}}]
\label{thm:BLSN-1}
Let $N > 1$ be an odd integer. Write $N - 1 = F R$ with $F > 1$,
$\gcd(F, R) = 1$, and $F$ completely factored into known primes.
Suppose for each prime $q \mid F$ there exists an integer $a_q$ with
\[
  a_q^{N-1} \;\equiv\; 1 \pmod N \qquad\text{and}\qquad
  \gcd(a_q^{(N-1)/q} - 1,\; N) \;=\; 1.
\]
If $F > \sqrt N$, then $N$ is prime. More generally, if $F > N^{1/3}$,
write $N = r F + s + 1$ with $0 \leq s < F$ and set
$\Delta = s^2 - 8 r$; if $\Delta$ is not a perfect square, then $N$ is
prime.
\end{theorem}

The $F > N^{1/3}$ threshold with the discriminant check
(\cite[Theorem~11]{BLS75}) is BLS's central improvement over the
classical $F > N^{1/2}$: instead of needing half the bit-size of $N$
to be factored, we need only a third.

\section{The three primality proofs}
\label{sec:proofs}

\subsection{Framework}

The three proofs in \S\S\ref{ssec:w2617}--\ref{ssec:w12391} invoke
Theorem~\ref{thm:BLSN-1} (BLS $N - 1$) with specific factorisations of
$N - 1 = 2(2^{p-1} - 1)/3$ drawn from the cyclotomic decomposition
\[
  2^{p-1} - 1 \;=\; \prod_{d \mid p - 1} \Phi_d(2),
\]
harvested partly from the Cunningham project tables~\cite{Cunningham}
and partly from direct computation. Each factor of each $\Phi_d(2)$
contributes to the BLS factored portion $F = F_-$ once its primality
is certified by APR-CL~\cite{APR,CohenLenstra}; Fermat-type witnesses
$a \in \{2, 3, 5, \ldots\}$ are then validated for each prime $q \mid
F$. The exact software toolchain and version numbers used to produce
the certificates are documented in Section~\ref{sec:comp}. The proofs
are unconditional and do not rely on the Chebyshev sufficiency
conjecture.

At each $p$, we also independently verify the Chua $N + 1$ condition
$\wom^{(W_p + 1)/2} \equiv -1 \pmod{W_p}$
(Proposition~\ref{prop:cond2}) as a sanity check on the
modular-arithmetic pipeline. This check is consistent with the BLS
proof (both confirm $W_p$ prime) but does not itself enter the BLS
argument.

APR-CL certification is applied uniformly to every prime $q$ of $F$,
regardless of size --- the smallest single-digit primes receive the
same unconditional APR-CL treatment as the largest multi-hundred-digit
factors. The ``exact margin'' column reports $(2F^3)_{\mathrm{bits}}
- N_{\mathrm{bits}}$, the integer bit-margin by which the BLS
hypothesis $2F^3 > N$ is satisfied.

\begin{center}
\begin{tabular}{rrrrrrr}
\toprule
$p$ & $W_p$ digits & div.\ of $p{-}1$ & primes in $F$ & exact margin (bits) & wall time \\
\midrule
$ 2617$ &  $788$ & $16$ &  $22$ &  $47$ & ${<}\,10$\,s (Ryzen 5 3600) \\
$10501$ & $3161$ & $48$ & $103$ & $3262$ & $98$\,s (1 core, 3.6\,GHz) \\
$12391$ & $3730$ & $32$ & $ 61$ & $2861$ & $82$\,s (1 core, 3.6\,GHz) \\
\bottomrule
\end{tabular}
\end{center}

For each proof, Condition~(II) is independently verified
($\wom^{(N+1)/2} \equiv -1 \pmod N$) as a sanity check; full data in
\texttt{bls\_certificate\_w\{2617,10501,12391\}.json}
(Section~\ref{sec:comp}).

\subsection{$W_{2617}$}
\label{ssec:w2617}

\begin{theorem}
\label{thm:w2617}
The Wagstaff number $W_{2617}$ is prime.
\end{theorem}

\begin{proof}
Set $N = W_{2617}$, a $788$-digit number. We apply
Theorem~\ref{thm:BLSN-1} to the $N - 1$ side; the factored portion
$F$ will be shown to exceed $N^{1/3}$.

The exponent satisfies $p - 1 = 2616 = 2^3 \cdot 3 \cdot 109$, with
$16$ divisors. The cyclotomic decomposition
\[
  2^{p-1} - 1 \;=\; \prod_{d \mid p-1} \Phi_d(2)
\]
gives $16$ factors, of which $12$ (those with $d \leq 654$) are fully
factored by direct computation. Their cumulative product contributes
$22$ prime factors to $F$ (a factor of $(N-1)/2$), and $2 F^3 > N$ by
a $47$-bit margin.

For each prime $q \mid F$ (the $22$ primes constituting the factored
portion), the witness $a_q \in \{2, 3\}$ is tested:
$a_q^{N-1} \equiv 1 \pmod N$ and $\gcd(a_q^{(N-1)/q} - 1, N) = 1$.
Every witness succeeds (full data in
\texttt{bls\_certificate\_w2617.json}). Writing $N = r F + s + 1$
with $0 \leq s < F$ and $\Delta = s^2 - 8 r$, the discriminant
$\Delta$ is not a perfect square.

By Theorem~\ref{thm:BLSN-1}, $N = W_{2617}$ is prime.
\end{proof}

\begin{remark}
As a sanity check, the Chua $N + 1$ condition $\wom^{(N+1)/2} \equiv
-1 \pmod N$ is verified independently: the modular exponentiation in
$\Z[\sqrttwo]/(N)$ takes under $1$ second with gmpy2.
\end{remark}

\subsection{$W_{10501}$}
\label{ssec:w10501}

\begin{theorem}
\label{thm:w10501}
The Wagstaff number $W_{10501}$ is prime.
\end{theorem}

\begin{proof}
Set $N = W_{10501}$, a $3161$-digit number. The exponent $p - 1 =
10500 = 2^2 \cdot 3 \cdot 5^3 \cdot 7$ has $48$ divisors. Of the $48$
cyclotomic factors $\Phi_d(2)$ of $2^{p-1} - 1$, $43$ are fully
factored --- $38$ by direct computation (using SymPy
\texttt{factorint} with Pollard $\rho$, $p - 1$, and ECM) and $5$
retrieved from the Cunningham project tables~\cite{Cunningham}. This
contributes $103$ primes to the BLS factored portion $F$, each
certified prime by APR-CL~\cite{APR,CohenLenstra}.

The factored portion satisfies $2 F^3 > N$ by a margin of $3262$
bits, substantially exceeding the BLS threshold. Fermat-type
witnesses $a_q \in \{2, 3, 5\}$ are validated for each of the $103$
primes $q \mid F$; every witness succeeds. The discriminant
$\Delta = s^2 - 8r$ (with $N = rF + s + 1$, $0 \leq s < F$) is
computed and verified non-square. Total wall time: $98$\,s on one
core of a 3.6\,GHz machine; full data in
\texttt{bls\_certificate\_w10501.json}.

By Theorem~\ref{thm:BLSN-1}, $N = W_{10501}$ is prime.
\end{proof}

\subsection{$W_{12391}$}
\label{ssec:w12391}

\begin{theorem}
\label{thm:w12391}
The Wagstaff number $W_{12391}$ is prime.
\end{theorem}

\begin{proof}
Set $N = W_{12391}$, a $3730$-digit number. The exponent $p - 1 =
12390 = 2 \cdot 3 \cdot 5 \cdot 7 \cdot 59$ has $32$ divisors. Of the
$32$ cyclotomic factors $\Phi_d(2)$, $27$ are fully factored: $19$
by direct computation, $7$ retrieved from the Cunningham
tables~\cite{Cunningham} and public factor databases (every such
factor re-certified by APR-CL~\cite{APR,CohenLenstra} before use),
and $1$ itself prime. This contributes $61$ primes in $F_-$, the
largest being a $371$-digit prime cofactor of $\Phi_{2065}(2)$. Each
of the $61$ primes receives independent APR-CL certification.

The factored portion $F$ satisfies $2 F^3 > N$ by a $2861$-bit
margin. Fermat-type witnesses $a_q \in \{2, 3\}$ are validated for
each of the $61$ primes; every witness succeeds. The discriminant
$\Delta = s^2 - 8r$ is non-square. Total wall time: $82$\,s on one
core of a 3.6\,GHz machine (less than $W_{10501}$ despite the larger
$N$, because $p - 1$ has fewer but larger prime factors: less
divisor enumeration, although each factorisation task is heavier).
Full data in \texttt{bls\_certificate\_w12391.json}.

By Theorem~\ref{thm:BLSN-1}, $N = W_{12391}$ is prime.
\end{proof}

\subsection{Limits of the method}

\begin{remark}[BLS ceiling]
\label{rem:bls-ceiling}
The BLS $N - 1$ approach requires $p - 1$ to be sufficiently smooth
that enough cyclotomic factors $\Phi_d(2)$ with $d \mid p - 1$ have
known factorisations. Among the $36$ known Wagstaff primes and
probable primes with $p > 2617$, the exponents $p = 10501$ and
$p = 12391$ are the only ones meeting this threshold in the current
Cunningham tables. The next candidate by smoothness, $p = 14479$
(with $p - 1 = 2 \cdot 3 \cdot 19 \cdot 127$), falls several thousand
bits short; the main missing piece is a partial factorisation of
$\Phi_{14478}(2)$. Extensions of the Cunningham
tables~\cite{Cunningham} --- or direct ECM on the divisors
$\Phi_d(2)$ for $d \mid p - 1$ with the target exponents $p$ ---
would push the feasibility frontier. The cofactor data recorded in
the JSON certificates (Section~\ref{sec:comp}) identifies exactly
which $\Phi_d(2)$ remain only partially factored for $p = 10501$ and
$p = 12391$; these are natural targets for further ECM work.
\end{remark}

\section{Computational details}
\label{sec:comp}

\subsection{Software toolchain}

All computational components in this paper are reproducible from the
scripts in the supplementary repository using free, open-source
software.

\begin{itemize}
\item \emph{Integer factorisation} (cyclotomic factors $\Phi_d(2)$):
  SymPy~$\geq$~1.12 \texttt{factorint} (combining trial division,
  Pollard $\rho$, Pollard $p-1$, and ECM internally), augmented for
  cyclotomic factors by trial division along the arithmetic
  progression $r \equiv 1 \pmod{d}$ implied by primitive-divisor
  theory.

\item \emph{Pre-known large factorisations.} The Cunningham project
  tables~\cite{Cunningham} for cyclotomic factors $\Phi_d(2)$ with
  $d \leq 1200$, and FactorDB as a discovery aid for two larger terms
  ($\Phi_{2065}(2)$ and $\Phi_{2478}(2)$ in Theorem~\ref{thm:w12391}).
  Every retrieved factor is independently re-verified for primality
  by APR-CL before use; FactorDB plays no evidentiary role.

\item \emph{Primality of cofactors} (all factors used in BLS
  certificates): the APR-CL implementation in PARI/GP~$\geq$~2.15
  (\texttt{isprime(n, 2)}), which forces the APR-CL
  (Adleman--Pomerance--Rumely / Cohen--Lenstra) path and is
  unconditionally deterministic~\cite{APR,CohenLenstra}. The largest
  cofactor is a $371$-digit factor of $\Phi_{2065}(2)$
  (Theorem~\ref{thm:w12391}); the smallest are single-digit primes,
  but all receive APR-CL certification for uniformity.

\item \emph{Fast modular arithmetic} (the modular exponentiations
  $a^{N-1} \bmod N$ and $\wom^{(N+1)/2} \bmod N$ for $N = W_{10501},
  W_{12391}$): Python's built-in big-int \texttt{pow}, with GMP-backed
  acceleration via gmpy2~$\geq$~2.1 when installed.
\end{itemize}

The driver scripts
\texttt{bls\_n\_minus\_1\_w\{2617,10501,12391\}.py} in the
supplementary repository emit JSON certificates
(\texttt{bls\_certificate\_w*.json}) recording every cyclotomic
factor, prime decomposition, BLS witness, and primality method used,
along with full execution logs. A master verifier
\texttt{verify\_bls.py} re-reads each certificate from disk, treats
every claim in it as untrusted input, rebuilds $F$ from $v_q(N-1)$,
re-runs APR-CL on every prime, re-checks the witnesses and the
discriminant, and re-runs Condition~(II); any disagreement exits
non-zero.

\subsection{Supplementary material}

All data and scripts are archived on Zenodo (DOI forthcoming upon
publication), with a development mirror on GitHub at
\url{https://github.com/dolonet/wagstaff-bls-primality}.

The archive contains, for each $p \in \{2617, 10501, 12391\}$:

\begin{enumerate}
\item the BLS primality certificate (JSON), recording every
  cyclotomic factor $\Phi_d(2)$, its prime decomposition, the BLS
  witness $a$ for each prime $q \mid F$, and the APR-CL timings for
  each cofactor;
\item the driver script \texttt{bls\_n\_minus\_1\_w<p>.py}
  reproducing the certificate end-to-end;
\item the full APR-CL primality log for each cofactor, reproducible
  via \texttt{gp -q isprime.gp};
\item the independent Condition~(II) verification log.
\end{enumerate}

The archive is deterministic and version-pinned; running the driver
scripts on a clean machine reproduces every numeric result in this
paper. The master verifier \texttt{verify\_bls.py} reads each
certificate and re-derives $F$, the margin, witnesses, discriminant,
and Condition~(II) from scratch.

\section{Relation to the Chebyshev sufficiency conjecture}

The content of Sections~1--\ref{sec:comp} is independent of the
Chebyshev sufficiency conjecture --- the question of whether the
single congruence $\wom^{(W_p + 1)/2} \equiv -1 \pmod{W_p}$ alone
suffices to decide primality of $W_p$. That conjecture is necessary
(Proposition~\ref{prop:cond2}) and conjecturally sufficient; its
current structural status --- including a reduction to the Pell
Primitive Pair Conjecture (PPPC) --- is the subject of the companion
paper~\cite{PaperB}. A further companion paper~\cite{PaperD3} gives a
universal unconditional Frobenius primality test in $\Z[\sqrt 3]$ for
the broader family $(b^p \pm 1)/(b \pm 1)$; the ring and the base
there are different from the ones used in this paper
($\Z[\sqrttwo]$ and $\wom = 3 + 2\sqrttwo = (1 + \sqrttwo)^2$ here
vs. $\Z[\sqrt 3]$ and $\omega_3 = (2 + \sqrt 3)^2 = 7 + 4\sqrt 3$
there).

\appendix

\section{Reference data}
\label{sec:appendix}

\subsection{The 36 known Wagstaff primes and probable primes}

Proved prime ($p \leq 42737$, by ECPP~\cite{AtkinMorain} or by this
paper), excluding the trivial case $p = 3$ for which $W_3 = 3$:
\[
  p \in \{5, 7, 11, 13, 17, 19, 23, 31, 43, 61, 79, 101, 127, 167,
  191, 199, 313, 347, 701, 1709,
\]
\[
  2617, 3539, 5807, 10501, 10691, 11279, 12391, 14479, 42737\}
\]
--- $29$ exponents. Probable prime ($p > 42737$, not yet proved but
consistent with every applied pseudoprime test):
\[
  p \in \{83339, 95369, 117239, 127031, 138937, 141079, 267017\}
\]
--- $7$ exponents. The ``$36$ known'' count in the header includes both
groups; adding $p = 3$ would make $37$.

\subsection{Cyclotomic decomposition of $W_p \pm 1$}

For odd prime $p$ and $N = W_p$:
\[
  N - 1 \;=\; \frac{2^p - 2}{3} \;=\; \frac{2(2^{p-1} - 1)}{3},
  \qquad N + 1 \;=\; 4\, W_{p-2}.
\]
On the $N - 1$ side the cyclotomic decomposition
\[
  2^{p-1} - 1 \;=\; \prod_{d \mid p-1} \Phi_d(2)
\]
harvests factored portions from divisors $d$ of $p - 1$; each
$\Phi_d(2)$ is a positive integer whose prime factors lie in an
arithmetic progression $r \equiv 1 \pmod d$. This is the source of
the BLS factored portions in
\S\S\ref{ssec:w2617}--\ref{ssec:w12391}.

On the $N + 1$ side the identity $W_p + 1 = 4 W_{p-2}$
(Lemma~\ref{lem:fundamental}) gives a single cofactor $W_{p-2}$ whose
factorisation is harder in general but tractable for small $p$;
Theorem~\ref{thm:Nplus1} is phrased in terms of this decomposition.

\begin{thebibliography}{99}

\bibitem{APR}
L.~M. Adleman, C.~Pomerance, R.~S. Rumely,
\emph{On distinguishing prime numbers from composite numbers},
Ann. of Math. (2) \textbf{117} (1983), 173--206.

\bibitem{AtkinMorain}
A.~O.~L. Atkin, F.~Morain,
\emph{Elliptic curves and primality proving},
Math. Comp. \textbf{61} (1993), no. 203, 29--68.

\bibitem{BLS75}
J.~Brillhart, D.~H. Lehmer, J.~L. Selfridge,
\emph{New primality criteria and factorizations of $2^m \pm 1$},
Math. Comp. \textbf{29} (1975), no. 130, 620--647.

\bibitem{Cunningham}
J.~Brillhart, D.~H. Lehmer, J.~L. Selfridge, B.~Tuckerman, S.~S. Wagstaff, Jr.,
\emph{Factorizations of $b^n \pm 1$, $b = 2, 3, 5, 6, 7, 10, 11, 12$ Up to
High Powers}, 3rd ed., Contemporary Mathematics, vol.~22,
American Mathematical Society, 2002. Online updates at
\url{https://homes.cerias.purdue.edu/~ssw/cun/}.

\bibitem{Chua}
K.~S. Chua,
\emph{Chebyshev polynomials and higher order Lucas--Lehmer algorithm},
preprint arXiv:2010.02677, 2020.

\bibitem{CohenLenstra}
H.~Cohen, A.~K. Lenstra,
\emph{Implementation of a new primality test},
Math. Comp. \textbf{48} (1987), no. 177, 103--121.

\bibitem{Lehmer}
D.~H. Lehmer,
\emph{On Lucas's test for the primality of Mersenne's numbers},
J. London Math. Soc. \textbf{10} (1935), no. 3, 162--165.

\bibitem{Williams}
H.~C. Williams,
\emph{\'Edouard Lucas and Primality Testing},
CMS Series of Monographs and Advanced Texts \textbf{22},
Wiley-Interscience, 1998.

\bibitem{PaperB}
A.~Dolotov,
\emph{Reduction of the Wagstaff Chebyshev sufficiency conjecture to
the Pell Primitive Pair Conjecture},
companion paper (in preparation).

\bibitem{PaperD3}
A.~Dolotov,
\emph{A deterministic Frobenius primality test in $\Z[\sqrt 3]$ for
the family $(b^p \pm 1)/(b \pm 1)$},
companion paper (in preparation).

\end{thebibliography}

\end{document}
